\documentclass[10pt,a4paper]{article}
\usepackage[T1]{fontenc}
\usepackage[margin=2cm]{geometry}
\usepackage{isabelle,isabellesym}
\usepackage{amsmath,amssymb}
\usepackage{pdfsetup}

\urlstyle{rm}
\isabellestyle{it}

\begin{document}

\title{Sylvester's Catalecticant}
\author{Carles Mar\'in}
\maketitle

\begin{abstract}
Sylvester's theorem on the catalecticant~\cite{Sylvester1851} characterizes when a binary form
$f=\sum_i a_i\binom{n}{i}x^{n-i}y^i$ is a sum of $r$ $n$-th powers of linear forms $\alpha_k x+\beta_k y$: this happens
(for pairwise non-proportional $(\alpha_k:\beta_k)$) exactly when there is a nonzero form
$h=\sum_l c_l x^{r-l}y^l$ vanishing at every $(\alpha_k,\beta_k)$ whose coefficients annihilate the Hankel
(catalecticant) matrix of $f$: $\sum_l a_{j+l}c_l=0$ for $j+r\le n$. We formalize the necessity over any commutative
ring, and the sufficiency over any field, first in the affine case and then with a point at infinity.
\end{abstract}

\paragraph{Acknowledgement.} The formalization and the text of this document were prepared with the
assistance of generative AI systems, and checked by the author.

\tableofcontents

\parindent 0pt\parskip 0.5ex

\input{session}

\bibliographystyle{abbrv}
\bibliography{root}

\end{document}
